\section{Introduction}
Gradient Descent-based algorithms such as  Stochastic Gradient Descent (SGD) and its variations, are a fundamental tool in training neural networks, and are therefore the subject of intense theoretical scrutiny. Recent years have witnessed significant advancements in understanding their dynamics and the learning mechanisms of neural nets. Significant progress has been made, in particular, in the case of two-layer networks thanks, in part, to the so-called mean-field analysis \citep{mei2018mean,chizat2018global,rotskoff2022trainability,sirignano2020mean}). A large part of the theoretical approach focused on \emph{one-pass} optimization algorithms, where \emph{each iteration involves a new fresh batch of data}. In particular, in the mathematically solvable case of  high-dimensional synthetic Gaussian data, and a low dimensional a multi-index target function model, the class of functions efficiently learned by these \textit{one-pass} methods has been thoroughly analyzed in a series of recent works, and have been shown to be limited by the so-called information exponent \citep{arous2021online} and leap exponent \citep{abbe2022merged,abbe2023sgd}. These analyses have sparked many follow-up theoretical works over the last few months, see, e.g. \cite{damian2022neural,damian2023smoothing,dandi2023how,bietti2023learning,ba2023learning,moniri2023theory,mousavihosseini2023gradient,zweig2023symmetric}, leading to a picture where SGD dynamics is governed by the information exponent, that is, the order of the first non-zero Hermite coefficient of the target.

A common point between all these works, however, is that they considered the idealized situation where one process data one sample at a time, all of them independently identically distributed, without any repetition. This is, however, not a realistic situation in practice. Indeed, most datasets contains similar data-points, so that repetition occurs (see e.g. for repetition and duplicate in CIFAR \citep{barz2020do}) Additionally, it is common in machine learning to  repeatedly go through the same mini-batch of data multiple time over epoch. Finally, many gradient algorithms, often used in machine learning in Lookahead optimizers \citep{zhang2019lookahead}, Extragradient methods \citep{korpelevi1976extragradient}, or Sharpness Aware Minimization (SAM) \citep{foret2021sharpness}, are \emph{explicitly} few gradient steps over the exact same data-point. It is thus natural to wonder if using such extra-gradient algorithms, or simply reusing many time some data, would change the picture with respect to the current consensus.


There are good reason to believe this to be the case. Indeed, it has been observed very recently \citep{dandi2024benefits} that iterating twice over large ($O(d)$) datasets was enough to alter the picture reached in 
\cite{arous2021online,abbe2022merged,abbe2023sgd}. Indeed \cite{dandi2024benefits} showed that their exist functions that can be learned over few (just two) batch repetition, while they require a much larger (i.e. polynomial in $d$) number of steps otherwise.
%gradient descent when repeating over data batches \emph{surpasses} (that scale linearly with the input dimension) the limitations imposed by the information and leap exponents, achieving a positive correlation with the target function for a much broader class than the staircase functions \cite{abbe2021staircase}. 
While the set of function considered in this work was limited, the conclusion is indeed surprising, thus motivating the following question: 
\begin{center}
\textit{
Can data repetition increase the efficiency of SGD when learning any multi-index functions?
% Is there a fundamental gap of processing data online with respect to repeating larger batches? \\
% Are single-pass algorithms intrinsically sub-efficient with respect to the multiple-pass one?
} 
\end{center}

% {\bf the point is that the SGD ideal situation is not realistic and very fragile, and Dandi et al already show tojs in a very specific setting (O(1)) easy function, two iterations and O(D) batched reusined) Our point is to study what happens in  more generic siutation bla bla blabl argue for our setting... (i) proxy ofr realistic sutaiton where data or correlated data are reused, and (ii) in pratice people are foten using SAM, look head, federated learning...}

We establish a positive answer to the above question by the analysis of more realistic optimization schemes with respect to the standard One-Pass SGD. Indeed, the hardness exponents developed in the seminal works \cite{abbe2021staircase,arous2021online} are substantially bound to the idealized training scenario that considers i.i.d data. Slight modifications to this training scenario toward a more realistic setting starkly change the global picture. Such slight modifications include: a) non-vanishing correlations between different data points, that is a natural requirement for real datasets; b) processing the same data point multiple times in the optimization routine, i.e., a standard step in any algorithmic procedure when looping over different epochs. 

We consider a simple yet analytically tractable single-pass training scheme: processing one Gaussian sample at a time, but \emph{repeating the gradient step twice}. This minor adjustment, common in Lookahead optimizers \citep{zhang2019lookahead}, Extragradient methods \citep{korpelevi1976extragradient}, and Sharpness-Aware Minimization (SAM) \citep{foret2021sharpness}, significantly boosts efficiency. In fact, we will see that this approach matches the performance of our best algorithms in most cases.

We show in particular that most multi-index function are learned with total sample complexity either $O(d)$ or $O(d \log d)$ if they are related to the presence of a symmetry, without any need for preprocessing of the data that is performed in various works (See e.g., \cite{mondelli2018fundamental,luo2019optimal,maillard2020phase,chen2020learning}).  Our results thus demonstrate that shallow neural networks turns out to be significantly more powerful than previously believed \citep{arous2021online,abbe2022merged,abbe2023sgd,arnaboldi2023high} in identifying relevant feature in the data. In fact, they are shown in many (but not all) cases to saturate the bounds predicted by statistical queries (SQ)  \citep{damian2024computational} (rather than the more limited correlation statistical queries (CSQ) bounds that were characteristic of the former consensus \citep{arous2021online,abbe2023sgd}). Equivalently, it roughly  mirrors the prediction of the approximate message passing analysis in \citep{troiani2024fundamental}. Our conclusions follow from the rigorous analysis of the overlaps between the trained weight vectors with the low-dimensional relevant subspace that draws inspiration from the pivotal works of \cite{saad1995line,coolen2000dynamics,coolen2000line,arous2021online}.

Concurrently with the present work, \citep{lee2024neural} also studies the properties of SGD under data repetition in neural networks. We discuss the differences between our works in Section~\ref{subsec:single_index_discussion}.
